\section{Introduction}

Tool-using language models are often evaluated as if each query required a single decision. That assumption is too weak for many agent settings. A user may simultaneously ask for news and market impact, or for travel planning and exchange-rate information. In such cases the model must emit multiple tools in sequence, not merely rank a single option.

This paper studies that multi-tool regime. Our starting point is a simple failure mode: a stationary attention intervention that boosts one facet at every decoding step tends to keep boosting the same facet after the first tool has already been emitted. On a multi-tool query, this is the wrong inductive bias. The model needs a mechanism that makes the second step less like the first one.

The intervention we evaluate is an ontology-guided query bias defined on the same per-head basis used by our key-side steering experiments. Its practical implementation is deliberately simple: at each layer, we project the query onto the ontology subspace and add a signed multiple of that projection. Negative coefficients suppress ontology-aligned query components and act as a coverage bias. The hypothesis is direct: for multi-tool prompts, this suppression should improve second-tool recovery, whereas a stationary key-side bias on the same basis should either do nothing or over-concentrate on the first facet.

The evidence supports that hypothesis, but only in a narrow sense. On MetaTool Subtask4, the Q-side bias is consistently positive on both Qwen and Llama, while the stationary K-side bias is negative on both and catastrophic on Llama. The gain is modest in absolute size, so the paper cannot be sold as a broad headline accuracy breakthrough. Its value lies elsewhere: the sign pattern is stable, the null controls are decisive, and the basis-specificity story is stronger than the raw gain.

The paper makes three claims. First, query-space coverage bias is the productive direction for multi-tool selection on the current benchmark. Second, the ontology basis is load-bearing: replacing it with random or feature-shuffled controls erases the gain and destroys K-side stability. Third, the local attention perturbation analysis is not vacuous; its empirical ratios are tiny on both Qwen and Llama, which is consistent with using the bound as a diagnostic language for these interventions. The paper does \emph{not} claim a clean full head-to-head win against SEKA or AdaSEKA on a matched benchmark. Its claim is narrower: once the task requires sequential recovery of more than one tool, the intervention site and sign become more important than the generic fact that the method is low-rank.
